\documentclass[preprint,12pt,authoryear]{elsarticle}
\usepackage{amsmath,amssymb,booktabs,graphicx,tabularx,longtable,float}
\usepackage{placeins}
\usepackage[T1]{fontenc}
\usepackage{lmodern}
\usepackage[utf8]{inputenc}
\usepackage[hidelinks]{hyperref}

\begin{document}

\begin{frontmatter}
\title{Supplementary Material for ``Centered Natural-Parameter Group Selection for Posterior-Score Contrast Support in von Mises--Fisher Mixtures''}
\author[aff1]{Hyunsoo Shin}
\author[aff1]{Byungtae Seo\corref{cor1}}
\cortext[cor1]{Corresponding author.}
\address[aff1]{Department of Statistics, Sungkyunkwan University, Seoul 03063, Republic of Korea}
\ead{seobt@skku.edu}
\end{frontmatter}

\renewcommand{\thetable}{S\arabic{table}}
\renewcommand{\thefigure}{S\arabic{figure}}
\renewcommand{\thesection}{S\arabic{section}}
\renewcommand{\theequation}{S\arabic{equation}}

\section{E-ACGL extension and path sensitivity}
\label{sec:s-eacgl-path}

The main article defines E-CGL as the primary estimator and E-ACGL as its
adaptive extension. This Supplementary Material gives the complete E-ACGL path specification, audit, and numerical
results. Notation and abbreviations not redefined here follow the main article.
Let $E^{\mathrm D}$ denote the common dense pilot used in the main
article. For completeness, with $c_{\cdot j}=H_KE_{\cdot j}$, its penalty is

\begin{equation}
\label{eq:supp-eacgl-penalty}
\begin{aligned}
\mathcal P_{\mathrm{ACGL}}(E)
&=\sum_{j=1}^d
w_j\lVert H_KE_{\cdot j}\rVert_2
=\sum_{j=1}^d
w_j\lVert c_{\cdot j}\rVert_2,\\
w_j^{\mathrm{raw}}
&=\left(\lVert H_KE_{\cdot j}^{\mathrm D}\rVert_2
+\epsilon\right)^{-\gamma},
\qquad
w_j
=\frac{w_j^{\mathrm{raw}}}
{\operatorname{median}_{1\leq h\leq d}w_h^{\mathrm{raw}}}.
\end{aligned}
\end{equation}

The weights are computed once from the common dense initializer and then held
fixed. For
$V=E^{(u)}-s\nabla f_t(E^{(u)})$, the weighted proximal step is

\begin{equation}
\label{eq:supp-eacgl-prox}
E_{\cdot j}^{(u+1)}
=J_KV_{\cdot j}
+\left(
1-\frac{s\lambda_\eta w_j}
{\lVert H_KV_{\cdot j}\rVert_2}
\right)_+H_KV_{\cdot j}.
\end{equation}

The E-step and guard logic are unchanged in form, but the criteria evaluated
by the guards are method-specific:

\[
Q_{\lambda_\eta}^{(\mathrm{ACGL})}
=Q-\lambda_\eta\sum_{j=1}^d
w_j\lVert H_KE_{\cdot j}\rVert_2,
\qquad
\mathcal L_{\lambda_\eta}^{(\mathrm{ACGL})}
=\ell-\lambda_\eta\sum_{j=1}^d
w_j\lVert H_KE_{\cdot j}\rVert_2.
\]

The reported analyses use $\gamma=1$ and $\epsilon=10^{-6}$; the weights are
not updated along the path. This construction applies the adaptive lasso and
adaptive group-lasso principles of \citet{zou2006adaptive} and
\citet{wang2008adaptivegroup}, respectively, to centered coordinate groups.
No oracle-property claim is made.

The final method-specific Karush--Kuhn--Tucker (KKT)-motivated path scale and
single path are

\begin{equation}
\label{eq:supp-eacgl-paths}
\begin{aligned}
\widetilde\lambda_{\max,\eta}^{\mathrm{ACGL}}
&=1.05\max_j
\frac{\lVert(H_KG_0)_{\cdot j}\rVert_2}{w_j},\\
\Lambda_{\mathrm{final}}
&=\{0\}\cup
\operatorname{Geom}
\left(10^{-4}\widetilde\lambda_{\max,\eta}^{\mathrm{ACGL}},
\widetilde\lambda_{\max,\eta}^{\mathrm{ACGL}};239\right),\\
\Lambda_{\mathrm{std}}^{\mathrm{hist}}
&=\{0\}\cup
\operatorname{Geom}
\left(10^{-2}\widetilde\lambda_{\max,\eta}^{\mathrm{ACGL}},
\widetilde\lambda_{\max,\eta}^{\mathrm{ACGL}};239\right).
\end{aligned}
\end{equation}

Here $E_0^{\mathrm{com}}$ is the matrix that repeats the pooled common natural
parameter, $\tau^0$ contains the responsibilities from the guarded zero-penalty
fit, and $G_0=\nabla f_{\tau^0}(E_0^{\mathrm{com}})$, as defined for the E-CGL
path scale in the main article. The final and
historical paths start from the same guarded zero-penalty fit obtained from
the common dense pilot. The final path processes its positive
tuning values in increasing order and uses the same support-deduplication,
target-preserving refit, explicit empty-support candidate, and BIC rule as
E-CGL. When several valid path fits generate the same support, the fit with
the largest observed log-likelihood initializes that support's refit.

Here $\operatorname{Geom}(a,b;q)$ denotes $q$ positive log-geometric values
from $a$ through $b$. The $10^{-2}$ path and the earlier two-segment audit that
adjoined values down to $10^{-6}$ are retained below as numerical provenance.
They motivated replacing the earlier lower endpoint by the single
$10^{-4}$ path, but they do not define the final estimator.

\section{Simulation design and complete numerical results}
\label{sec:s-simulation-results}

The primary simulation contains 12 cells formed by
$n\in\{300,1000\}$, target oracle Bayes error
$e_B\in\{0.025,0.05,0.10\}$, and equal or heterogeneous concentration.
Each cell has 100 paired repetitions and six common methods, giving 7,200
method-repetition rows. The recorded output contains no error rows and no
duplicate cell-repetition-method keys. Among the 7,200 primary-grid rows, one
dense shared-concentration fit in the equal-concentration $e_B=0.10$, $n=300$
cell reached its iteration limit. In the additional pure-concentration
diagnostic, four dense free-concentration pilots reached their iteration
limits. All selected E-series fits
passed the recorded majorization,
line-search acceptance, and path-endpoint checks. The largest oracle-error
calibration discrepancy was 0.172 percentage points (0.00172 on the probability
scale). The centered mean-direction group lasso (M-CGL) was fitted as a
directional companion in the four prespecified
$e_B=0.05$ cells, adding 400 valid method-repetition rows. It was not
subsequently fitted to the remaining eight primary cells.

\subsection{Oracle-error-calibrated data generation and randomization}

Let $q_{\mathrm{block}}=4$ denote the number of decision coordinates assigned to each of the
$K=4$ component blocks. For a decision coordinate in block $g$, define
\[
u_{kj}=\begin{cases}
1, & k=g,\\
-1/(K-1), & k\ne g.
\end{cases}
\]
For a candidate common-background norm $A$, the natural parameters are
generated as
\[
\eta_{kj}(A)=
\begin{cases}
A/\sqrt{q_C}, & j\in S_C,\\
\alpha_k(A)u_{kj}, & j\in S_D,\\
0, & j\in S_N,
\end{cases}
\qquad
\alpha_k(A)=\frac{\sqrt{\kappa_k^2-A^2}}{\lVert u_{k\cdot}\rVert_2}.
\]
Thus $\lVert\eta_k(A)\rVert_2=\kappa_k$, and
$\mu_k=\eta_k/\kappa_k$. Mixture proportions are $1/K$.

The value of $A$ is chosen by 18 bisection steps on
$[0,0.995\min_k\kappa_k]$. A common-random-number Monte Carlo sample of size
50,000 is used during bisection, followed by an independent validation sample
of size 200,000. The final achieved errors are reported in the main article.
For target $e_B$ and scenario index $h\in\{1,2\}$ (equal, heterogeneous), the
historical DGP calibration uses
\[
s_{\mathrm{cal}}
=271828182+100\,\operatorname{round}(1000e_B)+h.
\]
For the final rerun, let $c=1,\ldots,24$ denote the fixed row index in the
archived cell manifest and let $b=1,\ldots,100$ denote the repetition. The
production seeds are
\[
\begin{aligned}
s_{\mathrm{data},c,b}&=260731+100000c+b,\\
s_{\mathrm{test},c,b}&=260731+100000c+50000+b,\\
s_{\mathrm{init},c,b}&=260731+100000c+80000+b.
\end{aligned}
\]
These three seed ranges are disjoint across roles, cells, and repetitions.
The calibrated data-generating process (DGP) parameters are held fixed across
$n=300$ and $n=1000$ within each target-error and concentration combination.
Within a repetition, all methods are evaluated on the same sample and use the
same initialization seed. The likelihood-based mixture methods use the
prespecified 10-start initialization protocol appropriate to their
concentration parameterization.

Monte Carlo standard errors (MCSEs) were computed from the 100 replicates in each
cell. Across all methods and primary cells,
their maxima were 0.0094 for $F_{1,\eta}$, 0.0114 for ARI, and 0.101 for
$\operatorname{MSE}_{\eta^c}$. Restricted to E-CGL and E-ACGL, the
corresponding maxima were 0.0094,
0.0053, and 0.0375.

The primary selector is BIC after target-specific support refit. For E-CGL and
E-ACGL, every distinct path support and the explicit empty support are refitted
under an equality constraint that preserves the centered natural-parameter
target. The E-series nominal dimension is
\begin{equation}
\label{eq:supp-df-e}
\operatorname{df}_{E}(S)
=d+(K-1)m+(K-1)\mathbf 1(m>0),
\qquad m=|S|.
\end{equation}
It is a practical selection rule rather than an exact effective degrees of
freedom for a nonregular penalized mixture.

The Rossi-style entrywise mean-direction lasso (M-L) proposed by
\citet{rossi2022sparsevmf} uses its event-driven path with at most 240 path
points, including the zero-penalty fit, in the primary simulations. E-CGL and
E-ACGL use the zero-augmented 240-point
KKT-motivated path in the main article. The M-CGL companion uses the separate
240-point dense-score-proxy path defined in
Section~\ref{sec:s-target-diagnostic}. The E-series selector compares the distinct
refitted path supports with the all-common model
$S_\eta=\varnothing$. E-CGL and E-ACGL use the target-preserving
support-constrained refit and the nominal dimension in
\eqref{eq:supp-df-e}.
For the prototype-entry method M-L, an M-series mask with
$m_k\geq1$ active entries in component $k$ uses the operational dimension rule
\[
\operatorname{df}_{M,\mathrm{free}}
=(K-1)+K+\sum_{k=1}^K\max\{1,m_k-1\},
\]
for free component-specific concentrations, and
\[
\operatorname{df}_{M,\mathrm{shared}}
=(K-1)+1+\sum_{k=1}^K\max\{1,m_k-1\}
\]
for a shared concentration. The floor $\max\{1,m_k-1\}$ is the operational
BIC convention retained for this comparator. Thus these expressions are
practical method-specific dimensions rather than exact local dimensions;
they approximately account for mixing proportions, concentration parameters,
and unit-norm direction constraints.
They are not the centered M-CGL dimension, which is defined in
Section~\ref{sec:s-target-diagnostic}. Dense free- and shared-concentration
models have $Kd+K-1$ and $Kd$ degrees of freedom, respectively. E-ACGL weights
are fixed from the common dense component-specific-concentration pilot shared
by the two E-series methods, with $\gamma=1$ and $\epsilon=10^{-6}$; they are
not updated along the path.

Although M-L penalizes component entries, the reported $S_P$ metrics collapse
those entries to their coordinate union. They therefore assess
prototype-union recovery rather than recovery of the full component-by-coordinate
entry mask.

The finite vMF mixture likelihood can be unbounded as a component
concentration diverges.
Accordingly, dense initialization, the E-CGL/E-ACGL path, and their
support-constrained refits use the same finite parameter space,
\[
0\leq\kappa_k\leq\texttt{kappa\_cap}=10^6.
\]
Proximal proposals outside this space trigger step halving and are recomputed,
rather than projected. The support-constrained refit begins from a projected
and, if necessary, globally scaled source matrix, then step-halves each
unconstrained conditional proposal until the cap and likelihood guard hold. A
refit reaching at least
$0.99\,\texttt{kappa\_cap}$ is retained for diagnosis but is ineligible for
information-criterion selection. This bound defines the numerical problem and
does not act as a sparsity penalty.

\subsection{E-series numerical controls}

For each generalized M-step, the proximal-gradient loop stops when
\[
\frac{\lVert E^{(u+1)}-E^{(u)}\rVert_F}
{\max\{1,\lVert E^{(u)}\rVert_F\}}<10^{-8}
\]
or after 200 iterations in the numerical studies and 160 iterations in the
real-data analyses. The smooth majorization tolerance is $10^{-10}$,
and at most 40 backtracking attempts are allowed. The outer penalized
generalized expectation-maximization (GEM) loop uses relative objective
tolerance $10^{-7}$ and initially allows 100 iterations. Auxiliary- and
observed-objective decreases larger than
$10^{-8}$ mark the current path point as failed. Nonfinite likelihoods,
objectives, or natural parameters and concentration-cap violations are treated
likewise. A point that reaches its iteration limit is logged as unconverged
but remains in the provisional candidate-support family. An isolated failed
point is skipped, and the next tuning value is initialized from the last
accepted fit. Three consecutive numerical failures terminate the remaining
path without declaring convergence. If the penalized source of the provisional
refit-BIC winner is unconverged, the full path is rebuilt with 440 additional
iterations per point, only converged penalized sources are retained, and
support-constrained refitting and BIC selection are repeated. The path
also stops after storing a support with at most one active coordinate.

The support-constrained refit first allows 160 outer iterations and, if the
fit is finite, numerically valid, and unconverged, allows 440 additional
iterations in the recorded simulation and real-data protocols. Each conditional refit M-step uses the
limited-memory Broyden--Fletcher--Goldfarb--Shanno algorithm with bounds
(L-BFGS-B) \citep{byrd1995lbfgsb} and is limited to 80 iterations with
projected-gradient tolerance $10^{-6}$. The refit stops when
the absolute log-likelihood increment is at most
$10^{-7}+10^{-8}|\ell_{\mathrm{old}}|$. These values are numerical controls,
not statistical tuning parameters.
All reported fits use internally generated dense multistarts rather than the
optional externally supplied initialization, and every distinct support is
refitted; the optional refit shortlist is not used.

Table~\ref{tab:s-control-provenance} distinguishes settings used to obtain the
reported results from current software defaults. The numerical studies used
10 dense starts and the real-data analyses used 30. A later increase of the
package default for the additional refit budget from 440 to 840 iterations is
prospective and does not redefine any recorded fit.

\begin{table}[H]
\centering
\scriptsize
\setlength{\tabcolsep}{5pt}
\caption{Provenance of E-series numerical controls. The two recorded columns
define the analyses reported in the article; the last column records current
software defaults for settings controlled by \texttt{etaVmf}.}
\label{tab:s-control-provenance}
\begin{tabular}{lrrr}
\toprule
Control & Numerical studies & Real data & Current package default \\
\midrule
Path points (zero included) & 240 & 240 & analysis supplied \\
Relative positive-path floor & $10^{-4}$ & $10^{-4}$ & $10^{-4}$ \\
Initial outer penalized GEM iterations & 100 & 100 & 100 \\
Additional selected-source path iterations & 440 & 440 & 440 \\
Inner proximal iterations & 200 & 160 & 200 \\
Backtracking attempts & 40 & 40 & 40 \\
Initial refit iterations & 160 & 160 & 160 \\
Additional refit iterations & 440 & 440 & 840 \\
L-BFGS-B iterations per M-step & 80 & 80 & 80 \\
L-BFGS-B projected-gradient tolerance & $10^{-6}$ & $10^{-6}$ & $10^{-6}$ \\
Concentration cap & $10^6$ & $10^6$ & $10^6$ \\
\bottomrule
\end{tabular}
\end{table}

\subsection{Dense-pilot, selected-source, and path-stopping audits}
\label{sec:s-heuristic-audit}

Across the common 24-cell rerun, 2,396 of the 2,400 dataset-replicate dense
free-concentration pilots converged under the 200-iteration limit. Each pilot
was shared by E-CGL and E-ACGL, so this count is over dataset replicates rather
than 4,800 method-level fits. The four
remaining pilots occurred in the
pure-concentration cell at repetitions 4, 5, 39, and 56. Each terminal iterate
was continued under the same tolerance and concentration cap; convergence
required 8--60 additional iterations and increased the observed
log-likelihood by 0.30--7.72. Repeating the E-CGL and E-ACGL paths and all
support-constrained refits from these retry-converged pilots preserved all
eight selected supports and their $F_{1,\eta}$ values. The largest absolute
BIC change was 0.0264. If these eight sensitivity fits replaced the recorded
fits, the pure-concentration-cell mean ARI would change by less than 0.00013
and mean test NLL by less than 0.0012 for either method. The recorded
100-repetition summaries are therefore retained. The final implementation
continues a provisional winning pilot that reaches its iteration limit and,
if the continuation remains unconverged, falls back to the best converged
original start when one exists. Use of a finite unconverged pilot is flagged
and occurs only when no converged candidate is available.

The initial candidate archive contained 4,800 E-series method--replicate groups
and 328,703 support-constrained refits. Every selected refit was eligible and
converged, but this did not establish convergence of the penalized path fit
used to generate its support. We therefore reproduced all 4,800 selected
penalized sources from the final $10^{-4}$ paths. Thirty-seven selected sources
had reached the original 100-iteration limit, and 32 of their supports were
not generated by any converged point on the corresponding recorded path.

Only these 37 method--replicate paths were rerun with a 540-iteration allowance,
and the candidate family was reconstructed from converged path points only.
Every path point in the 37 reruns converged, and every selected
support-constrained refit passed the eligibility checks. Thirty-four selected
supports changed, including 32 changes in support size. The corrected archive
contains the same 8,300 sparse-method rows and 4,800 E-series
method--replicate groups, with 328,701 distinct-support refit records. It has no
error rows or duplicate keys, and every cell--method summary contains 100 valid
replications. The original archive is retained unchanged for provenance; all
additional and stress results in the article use the corrected archive.

Table~\ref{tab:s-selected-source-audit} localizes the 37 corrections. None
occurred in the primary 12-cell factorial grid. The pure-concentration change
affects one E-ACGL replication; the remaining changes occur in stress cells.

\begin{table}[t]
\centering
\small
\caption{Selected-source convergence corrections. ``Affected'' counts
archived winners whose penalized source reached the 100-iteration limit;
``support changed'' and ``$q$ changed'' compare the corrected converged-source
winner with the archived winner.}
\label{tab:s-selected-source-audit}
\resizebox{\textwidth}{!}{%
\begin{tabular}{llrrr}
\toprule
Cell & Method & Affected & Support changed & $q$ changed \\
\midrule
Pure concentration, $n=1000$ & E-ACGL & 1 & 1 & 1 \\
Moderately dense, $n=300$ & E-CGL & 2 & 2 & 2 \\
Moderately dense, $n=300$ & E-ACGL & 6 & 4 & 4 \\
Moderately dense, $n=1000$ & E-CGL & 3 & 3 & 3 \\
Moderately dense, $n=1000$ & E-ACGL & 1 & 1 & 0 \\
Strongly dense, $n=1000$ & E-CGL & 18 & 17 & 17 \\
Strongly dense, $n=1000$ & E-ACGL & 3 & 3 & 3 \\
High-dimensional, $n=300,d=500$ & E-CGL & 3 & 3 & 2 \\
\bottomrule
\end{tabular}}
\end{table}

The terminal finite path support has size zero or one in every simulation fit.
The support-size early stop was audited separately by reproducing 12
representative E-series paths under the initial 100-iteration allowance. The
audit includes both methods in the original three conditions,
the difficult heterogeneous $e_B=0.10$, $n=300$ cell, the $d=500>n$ cell, and
the final-path E-ACGL replicate with the smallest first-positive retention.
Each fit was continued from the first stored support of size at most one
through all 240 tuning values. Table~\ref{tab:s-full-path-audit} reports the
results. All 175 remaining points were accepted, no support with more than one
coordinate re-entered, and no new support appeared. Fourteen reproduced path
points reached their iteration limit and generated 11 supports not reproduced
by a converged path point, but none was the ultimately selected support.
Reading the continued fits at thresholds $10^{-7}$, $10^{-8}$, and $10^{-9}$
gave identical supports. These checks are targeted diagnostics rather than a
proof that support paths are nested; they are distinct from the exhaustive
4,800-winner source audit above.

Each recorded Classic3 full-data and split-specific E-series path contains all
240 configured points and ends with support size 2--4; early stopping therefore
truncates none. Of their 2,880 path points, 66 reached the original
100-iteration limit, but none generated a selected refit's source support.
Removing them leaves every selected Classic3 support unchanged. All selected
Classic3 refits are eligible and converged, and all path points pass the
recorded majorization check.

\begin{table}[t]
\centering
\scriptsize
\setlength{\tabcolsep}{3pt}
\caption{Targeted full-path audit after the support-size early stop. The
remaining points are the configured tuning values evaluated after the recorded
path had stopped. ``Unconv.-only'' counts supports generated by an
iteration-limit path point and by no converged point on the reproduced path;
none is the selected support.}
\label{tab:s-full-path-audit}
\resizebox{\textwidth}{!}{%
\begin{tabular}{llrrrrr}
\toprule
Cell & Method & Stored & Remaining & Unconv.-only & New & $q>1$ re-entry \\
\midrule
Heterogeneous, $e_B=.05$, $n=1000$ & E-CGL  & 227 & 13 & 0 & 0 & 0 \\
Heterogeneous, $e_B=.05$, $n=1000$ & E-ACGL & 226 & 14 & 0 & 0 & 0 \\
Pure concentration, $n=1000$       & E-CGL  & 236 & 4  & 4 & 0 & 0 \\
Pure concentration, $n=1000$       & E-ACGL & 234 & 6  & 0 & 0 & 0 \\
Moderately dense, $n=300$          & E-CGL  & 221 & 19 & 0 & 0 & 0 \\
Moderately dense, $n=300$          & E-ACGL & 214 & 26 & 0 & 0 & 0 \\
Heterogeneous, $e_B=.10$, $n=300$  & E-CGL  & 224 & 16 & 3 & 0 & 0 \\
Heterogeneous, $e_B=.10$, $n=300$  & E-ACGL & 222 & 18 & 1 & 0 & 0 \\
High-dimensional, $n=300,d=500$    & E-CGL  & 224 & 16 & 2 & 0 & 0 \\
High-dimensional, $n=300,d=500$    & E-ACGL & 219 & 21 & 1 & 0 & 0 \\
Crossed support, $n=400$, rep. 16   & E-CGL  & 229 & 11 & 0 & 0 & 0 \\
Crossed support, $n=400$, rep. 16   & E-ACGL & 229 & 11 & 0 & 0 & 0 \\
\bottomrule
\end{tabular}}
\end{table}

\subsection{Estimand-separation and stress data-generating models}

The pure-concentration model has $d=200$ and a common unit direction supported
on the first 16 coordinates with alternating signs. Its concentrations are
\[
\kappa=(10,30,80,200),
\]
so $|S_\mu|=0$ and $|S_\eta|=16$. The shared-natural-parameter model uses the
construction above with
\[
(q_C,q_D,q_N)=(80,20,100),
\qquad
\kappa=(30,40,50,60),
\]
and calibrated common norm $A=21.0118$ for target $e_B=0.05$.

The crossed-support diagnostic has $d=24$ and
$\kappa=(30,40,50,60)$. Four $\eta$-only coordinates use a common directional
loading of 0.20, four $\mu$-only coordinates use a common natural-parameter
loading of 8, eight coordinates complete the component row norms with
component-specific signed contrasts, and eight coordinates are null. Thus
$|S_\mu|=|S_\eta|=12$, with four coordinates unique to each target and eight
in their intersection.

The moderately and strongly dense controls use $d=200$, common norm
$A=27.7349$, and
\[
\kappa=(30,40,50,60),
\]
with
\[
(q_C,q_D,q_N)=(4,80,116)
\quad\text{or}\quad
(4,160,36),
\]
respectively. Both were calibrated to an oracle Bayes error of approximately
0.10. The
high-dimensional cell uses $n=300$, $d=500$,
$(q_C,q_D,q_N)=(10,40,450)$, $\kappa=(45,60,75,90)$, and calibrated
$A=28.5773$ for target $e_B=0.05$. The beta-min cell uses $n=1000$, $d=200$,
$(q_C,q_D,q_N)=(4,16,180)$, $\kappa=(30,40,50,60)$, and calibrated
$A=21.1746$. Within each four-coordinate decision block, one contrast has
unit weight and three have weight 0.25. These constructions and their
invariant checks are the recorded models used by the reported simulations.

\subsection{Partition, support, and parameter metrics}
\label{sec:support-metrics}

For the true partition $\mathcal U$ and fitted partition $\mathcal V$, the
reported NMI uses geometric-mean entropy normalization,
\[
\operatorname{NMI}(\mathcal U,\mathcal V)
=\frac{I(\mathcal U;\mathcal V)}
{\{H(\mathcal U)H(\mathcal V)\}^{1/2}}.
\]
It is set to one when the two partitions are identical and both empirical
entropies vanish, and to zero when exactly one empirical entropy vanishes.
ARI is computed in its standard chance-adjusted form.

For an estimated coordinate support $\widehat S$, the reported decomposition
is $\widehat q_C=|\widehat S\cap S_C|$,
$\widehat q_D=|\widehat S\cap S_D|$, and
$\widehat q_N=|\widehat S\cap S_N|$. The posterior-score contrast support
metrics are
\[
\operatorname{TPR}_\eta=\frac{\widehat q_D}{q_D},\qquad
\operatorname{FPR}_\eta=\frac{\widehat q_C+\widehat q_N}{q_C+q_N},\qquad
\operatorname{Precision}_\eta(\widehat S)
=\frac{\widehat q_D}{|\widehat S|},
\]
with precision set to zero for an empty support. When the denominator is
positive,
\[
F_{1,\eta}
=\frac{2\operatorname{TPR}_\eta\,
\operatorname{Precision}_\eta(\widehat S)}
{\operatorname{TPR}_\eta+\operatorname{Precision}_\eta(\widehat S)}.
\]
When the true support is nonempty and the estimated support is empty, TPR,
precision, and $F_{1,\eta}$ are set to zero. More generally, $F_{1,\eta}=0$
whenever
$\operatorname{TPR}_\eta+\operatorname{Precision}_\eta(\widehat S)=0$.
For a nonempty true support, exact recovery is recorded by
$\operatorname{Exact}_\eta(\widehat S)
=\mathbf 1\{\widehat S=S_\eta^\star\}$, and the exact-support rate is its
Monte Carlo mean. When the true support is empty, these
three quantities are undefined and displayed as dashes; exact-empty recovery
is reported instead. Target-specific prototype and directional evaluations
replace the subscript $\eta$ by $P$ and $\mu$, respectively.
Before parameter MSE is computed, all $K!$ component permutations
are examined and the permutation minimizing centered natural-parameter MSE is
used for $\mu$, $\kappa$, and centered-$\eta$ comparisons.

Using $(\eta^c)_{kj}=c_{kj}$ for a centered natural-parameter entry, write
$c_{kj}^{\star}=\eta_{kj}^{\star}
-K^{-1}\sum_h\eta_{hj}^{\star}$ and
$\widehat c_{kj}=\widehat\eta_{kj}
-K^{-1}\sum_h\widehat\eta_{hj}$. The parameter error reported throughout
the article and Supplementary Material is
\[
\operatorname{MSE}_{\eta^c}
=\frac{1}{Kd}\sum_{k=1}^K\sum_{j=1}^d
(\widehat c_{kj}-c_{kj}^{\star})^2.
\]
It is an error for centered natural-parameter contrasts, rather than for the
uncentered natural parameters.

Tables~\ref{tab:s42-support}--\ref{tab:s45-convergence} provide the
complete primary-grid, sample-size, oracle-benchmark, stress, and numerical
diagnostic results summarized in the main article.
The machine-readable companion file
\nolinkurl{supplementary_simulation_summary.csv} contains the method--cell
means and standard deviations used for the numerical tables, including the
final single-path E-ACGL summaries for the 12 primary and five stress cells and
the target-specific support metrics needed to reproduce Monte Carlo standard
errors.

\input{../simulations/tables/csda_supp_numerical_tables_260803.tex}

\clearpage

\subsection{Empty-support path sensitivity}

The E-series candidate family includes the all-common model even when the
finite path does not generate it. Under the historical $10^{-2}$-floor path,
we evaluated this candidate retrospectively for both E-series methods in 24
simulation cells, giving 4,800 method--repetition comparisons. It changed 84
selections, all in the three stress combinations in
Table~\ref{tab:s-empty-support}; primary factorial, estimand-separation, and
sample-size results were unchanged. These rows document the diagnostic that
motivated the wider candidate path and are not final performance results. On
the final single $10^{-4}$-floor path, neither E-CGL nor E-ACGL selected the
empty support in any of the 2,400 recorded method--repetition fits per method
in this common 24-cell rerun subset.

\begin{table}[H]
\centering
\small
\caption{Historical $10^{-2}$-floor audit of the explicit empty-support
candidate. Rows list the stress cells in which at least one selected model
changed; these counts are retained as path-development provenance.}
\label{tab:s-empty-support}
\begin{tabularx}{\textwidth}{@{}Xlrr@{}}
\toprule
Stress condition & Method & Comparisons & Empty-support wins \\
\midrule
Moderately dense, $n=300,d=200,q_D=80$ & E-CGL  & 100 & 81 \\
Moderately dense, $n=300,d=200,q_D=80$ & E-ACGL & 100 & 1 \\
High-dimensional, $n=300,d=500,q_D=40$ & E-CGL  & 100 & 2 \\
\bottomrule
\end{tabularx}
\end{table}

\section{Matched penalty-structure ablation}
\label{sec:s-penalty-ablation}

E-CL denotes the entrywise centered lasso with

\[
\mathcal P_{\mathrm{E-CL}}(E)
=\sum_{j=1}^d\sum_{k=1}^K|c_{kj}|,
\qquad c_{\cdot j}=H_KE_{\cdot j}.
\]

The fitted E-CL criterion subtracts
$\lambda_\eta\mathcal P_{\mathrm{E-CL}}(E)$.

E-CL, E-CGL, and E-ACGL were compared using the matched smooth
objective, paired samples and dense initializations, method-specific KKT
endpoints, 120-point paths, and target-preserving support-constrained refits.
Every distinct path support was refitted and compared by
$\operatorname{df}_E$-based BIC after refit. Each cell used
100 repetitions (Table~\ref{tab:s-penalty-ablation}). The quantity
$\widehat q_{\mathrm{part}}$ counts coordinates for which the penalized source
solution retained only a proper subset of the $K$ centered component entries;
it is evaluated before support refitting.

The well-separated S1 cell has $K=4$, $n=1000$, $d=200$, and
$q_D=16$. The difficult cell has $K=4$, $n=300$, $d=200$,
$(q_C,q_D,q_N)=(4,16,180)$, heterogeneous concentrations, and target
$e_B=0.10$.

\begin{table}[H]
\centering
\footnotesize
\setlength{\tabcolsep}{2.5pt}
\caption{Matched ablation of centered penalty structures. Entries are means
over 100 paired replications per cell.}
\label{tab:s-penalty-ablation}
\begin{tabular*}{\textwidth}{@{\extracolsep{\fill}}llrrrrrrrr@{}}
\toprule
Cell & Method & $\widehat q$ & $\widehat q_C$ & $\widehat q_D$ & $\widehat q_N$ & $\widehat q_{\mathrm{part}}$ & $F_{1,\eta}$ & ARI & $\operatorname{MSE}_{\eta^c}$ \\
\midrule
S1 & E-CL & 16.00 & 0.00 & 16.00 & 0.00 & 1.41 & 1.000 & 0.875 & 0.058 \\
S1 & E-CGL & 16.00 & 0.00 & 16.00 & 0.00 & 0.00 & 1.000 & 0.875 & 0.058 \\
S1 & E-ACGL & 16.00 & 0.00 & 16.00 & 0.00 & 0.00 & 1.000 & 0.875 & 0.058 \\
difficult $n=300$ cell & E-CL & 26.51 & 0.22 & 15.09 & 11.20 & 26.23 & 0.737 & 0.606 & 1.132 \\
difficult $n=300$ cell & E-CGL & 20.28 & 0.10 & 14.70 & 5.48 & 0.00 & 0.822 & 0.633 & 0.923 \\
difficult $n=300$ cell & E-ACGL & 15.32 & 0.00 & 14.85 & 0.47 & 0.00 & 0.948 & 0.706 & 0.413 \\
\bottomrule
\end{tabular*}
\end{table}

The well-separated S1 cell did not distinguish the three penalties by support
recovery. In the difficult cell, E-CGL selected 5.72 fewer noise coordinates
than E-CL and increased $F_{1,\eta}$ by 0.085; the paired Monte Carlo standard errors
were 0.82 and 0.010, respectively. E-ACGL selected 5.01 fewer noise
coordinates than E-CGL and increased $F_{1,\eta}$ by 0.126, with paired Monte Carlo
standard errors 0.53 and 0.009. Thus, coordinate grouping reduced entry
fragmentation and false selection in the difficult cell, while adaptive
weighting provided an additional improvement. The ablation isolates the
effect of fixed adaptive weighting and does not imply uniform superiority
over all regimes.

\section{Selector sensitivity}
\label{sec:s-selector-sensitivity}

A proximal-path sensitivity diagnostic compared BIC before and after
refitting. It used the main nominal dimension
$\operatorname{df}_A(S):=\operatorname{df}_E(S)$ in
\eqref{eq:supp-df-e},
and the active-only sensitivity dimension

\[
\operatorname{df}_B=Km+(K-1)\mathbf 1(m>0),
\qquad m=|S|.
\]

\begin{samepage}
For support selection, the EBIC sensitivity criterion
\citep{chen2008ebic} was implemented as
\[
\begin{aligned}
\operatorname{EBIC}_{\gamma}(S)
&=-2\ell\!\left(\widehat\Theta_S^{\mathrm{refit}}\right)
+\log(n)\operatorname{df}_A(S)\\
&\quad+2\gamma\log\binom{d}{|S|},
\qquad \gamma\in\{0.25,0.5,1\}.
\end{aligned}
\]
\end{samepage}
The two dimension rules are practical penalty slopes rather than exact
effective degrees of freedom. Within this five-replication audit,
$\gamma=0.25$ and $0.5$ selected identical supports within each method. For
E-CGL, both gave mean $\widehat q=17.2$, $\widehat q_D=13.8$,
$\widehat q_N=3.4$, $F_{1,\eta}=0.834$, and ARI $=0.631$; for E-ACGL, both
selected the same supports as BIC after refitting under
$\operatorname{df}_A$. Table~\ref{tab:s-selector-sensitivity} shows the most
difficult $n=300$ cell, displays $\gamma=1$ as the strongest EBIC setting,
and omits the duplicate intermediate rows. This five-replication diagnostic
is interpreted as a selector audit rather than a performance experiment.

\begin{table}[H]
\centering
\footnotesize
\setlength{\tabcolsep}{3pt}
\caption{Sensitivity to the information criterion in the difficult
heterogeneous cell. Entries are means over five diagnostic replications.}
\label{tab:s-selector-sensitivity}
\begin{tabular}{llrrrrr}
\toprule
Method & Rule & $\widehat q$ & $\widehat q_D$ & $\widehat q_N$ & $F_{1,\eta}$ & ARI \\
\midrule
E-CGL & BIC after, df-A & 19.0 & 14.2 & 4.8 & 0.814 & 0.655 \\
E-CGL & BIC before, df-A & 14.6 & 13.4 & 1.2 & 0.874 & 0.627 \\
E-CGL & BIC after, df-B & 14.8 & 13.6 & 1.2 & 0.882 & 0.616 \\
E-CGL & EBIC$_1$ after, df-A & 14.8 & 13.6 & 1.2 & 0.882 & 0.616 \\
E-ACGL & BIC after, df-A & 14.8 & 14.6 & 0.2 & 0.945 & 0.692 \\
E-ACGL & BIC before, df-A & 15.4 & 14.4 & 1.0 & 0.916 & 0.694 \\
E-ACGL & BIC after, df-B & 14.2 & 14.2 & 0.0 & 0.937 & 0.684 \\
E-ACGL & EBIC$_1$ after, df-A & 14.2 & 14.2 & 0.0 & 0.937 & 0.684 \\
\bottomrule
\end{tabular}
\end{table}

The stronger criteria removed false positives but also removed posterior-score
coordinates in the difficult cell. BIC after target-preserving
support-constrained refit with $\operatorname{df}_A$ is retained
as the primary practical rule; no exact effective-df claim is made.

Table~\ref{tab:s-path-coverage} separates two events that a single exact-support
recovery rate conflates. Path coverage asks whether the fitted path ever
contained the true support; conditional BIC success asks whether BIC selected
it when it was available. In the difficult heterogeneous trajectory, the
principal small-sample limitation was incomplete path coverage. The
$e_B=0.10$ selector gap was not monotone in $n$, although both path coverage
and the oracle-NLL gap improved markedly by $n=2000$.

\input{../simulations/tables/csda_section44_path_coverage_conditional_bic_260808.tex}

\subsection{Classic3 path sensitivity}

The original E-series analysis used a single 240-point path with relative
lower endpoint $10^{-2}$. E-CGL moved gradually from the zero-penalty fit, but
the adaptive weights caused the first positive E-ACGL point to move from
$q=2000$ to $q=442$. A two-segment audit that extended the candidate family
down to $10^{-6}$ showed that the dense E-ACGL selection was a
candidate-coverage issue rather than evidence that adaptive weighting intrinsically
requires a dense support.

To retain one path per method rather than a dataset-specific union, the final
analysis uses a method-specific single $L=240$ path from
$10^{-4}\widetilde\lambda_{\max}$ to $\widetilde\lambda_{\max}$ for each
E-series method. The data, dense starts,
fixed weights, target-preserving refit, explicit empty support, and BIC rule
are unchanged. Absolute penalty values remain method-specific because the
path scale contains the fixed weights.

For E-CGL, the final first positive value was 0.010361 and retained all 2,000
coordinates. The selected refit had $q=1420$, BIC $-37{,}798{,}212$, ARI
0.9761, and NMI 0.9539. For E-ACGL, the final first positive value was 0.196404
and retained 1,964 coordinates. The selected refit had $q=1438$, BIC
$-37{,}798{,}237$, ARI 0.9746, and NMI 0.9515. The earlier two-segment audit
selected $q=1421$ for E-CGL and $q=1435$ for E-ACGL; those values are retained
only as path-sensitivity provenance.

\begin{table}[H]
\centering
\scriptsize
\setlength{\tabcolsep}{2.5pt}
\caption{Candidate-path sensitivity of the E-series methods in Classic3.
The final rows use one 240-point path with relative lower endpoint $10^{-4}$;
the other rows record the preceding diagnostics.}
\label{tab:s-classic3-paths}
\resizebox{\textwidth}{!}{%
\begin{tabular}{lllrrrrr}
\toprule
Method & Candidate rule & Floor ratio & First $\lambda>0$ & First $q$ & Selected $\lambda$ & Selected $q$ & BIC \\
\midrule
E-CGL & historical single path & $10^{-2}$ & 1.036103 & 1947 & 3.1827 & 1421 & $-37{,}798{,}211$ \\
E-CGL & final single path & $10^{-4}$ & 0.010361 & 2000 & 3.1827 & 1420 & $-37{,}798{,}212$ \\
E-ACGL & historical single path & $10^{-2}$ & 19.6404 & 442 & 0 & 2000 & $-37{,}793{,}072$ \\
E-ACGL & final single path & $10^{-4}$ & 0.196404 & 1964 & 2.0024 & 1438 & $-37{,}798{,}237$ \\
\bottomrule
\end{tabular}%
}
\end{table}

\subsection{Simulation path sensitivity}

The standard-path numerical-study outputs were archived before applying the
final E-ACGL candidate rule. We screened 5,025
archived E-series replicate--method candidate rows across 26 cells. For each
row, the first-positive retention ratio was the support size at the first
positive standard-path value divided by its zero-penalty support size.
Both methods had 2,600 archived method--replicate groups with positive and
BIC-eligible candidates. In 175 E-CGL groups, the support-deduplicated
candidate table did not retain a separate zero-penalty representative because
that fitted support was represented by a positive tuning value. The retention
ratio was therefore computable for 2,425 E-CGL groups and all 2,600 E-ACGL
groups.
Table~\ref{tab:s-support-retention} shows that the E-CGL standard path moved
gradually from zero, whereas the adaptive weights often produced a much larger
first E-ACGL support jump. This screening is a candidate-coverage diagnostic;
it does not establish that an uncomputed near-zero candidate would have won
the refit-and-BIC comparison.

\begin{table}[H]
\centering
\small
\setlength{\tabcolsep}{4pt}
\caption{First-positive support retention under the recorded standard
E-series paths. Entries summarize the ratio of the first-positive support size
to the zero-penalty support size across method--replicate rows.}
\label{tab:s-support-retention}
\begin{tabular}{lrrrrrrrr}
\toprule
Method & Rows & Minimum & 1\% & 5\% & 25\% & Median & $<0.90$ & $<0.50$ \\
\midrule
E-CGL  & 2425 & 0.917 & 0.950 & 0.958 & 0.990 & 0.995 & 0 & 0 \\
E-ACGL & 2600 & 0.070 & 0.080 & 0.080 & 0.120 & 0.350 & 2369 & 1609 \\
\bottomrule
\end{tabular}
\end{table}

We then refitted both E-series methods on the augmented union in five targeted
cells: the heterogeneous $e_B=0.05$, $n=1000$ primary cell; the crossed-support
$n=400$ cell; the pure-concentration $n=1000$ cell; the equal-concentration
$e_B=0.05$, $n=2000$ cell; and the weak-beta-min $n=1000$ cell. Replicates
1--3 were used in each cell for both methods, giving 30 paired fits. All fits
were valid. E-CGL selected the standard segment in all 15 repetitions. E-ACGL
selected a candidate generated by the added segment in three repetitions, but
each had exactly the same support as the archived selection. Across the 30
comparisons, selected $q$, support, $F_{1,\eta}$, and ARI were unchanged; the
largest absolute BIC difference was $4.89\times10^{-4}$. The targeted audit
detected no change in the checked fits.

As a historical sensitivity analysis, E-ACGL was also rerun under the
two-segment rule in all 12 primary cells and five stress cells. The exact
selected support was unchanged in all 1,200 primary repetitions and in four
of the five stress cells. In the weak beta-min cell, an added candidate changed
the exact support in 19 repetitions: mean selected size changed from 15.80 to
15.99, $F_{1,\eta}$ from 0.994 to 0.998, exact recovery from 0.81 to 0.93, and
ARI from 0.869 to 0.870. Table~\ref{tab:s-eacgl-path-sensitivity} records this
audit without defining the final estimator or implying superiority of
adaptive weighting.

\clearpage
\input{../simulations/tables/csda_supp_eacgl_augmented_core_260809.tex}

The historical two-segment elapsed times include both segments, support
refits, and selection after reuse of the paired dense initialization. They are
descriptive run records and are not compared with the end-to-end
matched-runtime panel.

Both E-series methods were subsequently rerun on the final single
$10^{-4}$-floor path in all 12 primary cells, three estimand-separation cells,
four $e_B=0.05$ sample-size cells, and five stress cells. All 4,800 selected
refits under the support constraints were finite and converged. Their
penalized source fits were audited separately for convergence. The numerical
tables in this Supplementary Material use these final single-path results;
the preceding standard and two-segment results are retained only in the
path-audit tables.

The final hard-trajectory extension added the remaining heterogeneous
$e_B=0.10$ sample-size fits. Across all recorded final-path cells, E-CGL has
2,800 valid fits in 28 cells and E-ACGL has 2,600 valid fits in 26 cells; no
fit was nonfinite or nonconverged. The unequal totals arise because the
complete hard trajectory was run only for E-CGL.

\section{Ambient-dimension sensitivity}
\label{sec:s-ambient-dimension}

The dimension experiment retains $q_C=4$ and $q_D=16$ while varying the number
of noise coordinates. Table~\ref{tab:s-ambient-dimension} pools equal and heterogeneous
concentration cells and uses target $e_B=5\%$. At $d\in\{100,500\}$, each
concentration pattern has 50 repetitions, giving 100 after pooling; at $d=200$,
each pattern has 100 repetitions, giving 200 after pooling. These archived
sensitivity fits use 120-point standard paths.

\begin{table}[H]
\centering
\footnotesize
\setlength{\tabcolsep}{3pt}
\caption{Ambient-dimension sensitivity of E-CGL and E-ACGL. Entries are means
pooled across equal and heterogeneous concentration patterns under the
archived 120-point standard-path protocol.}
\label{tab:s-ambient-dimension}
\begin{tabular}{rrlrrrrr}
\toprule
$n$ & $d$ & Method & $\widehat q$ & $\widehat q_D$ & $\widehat q_N$ & $F_{1,\eta}$ & ARI \\
\midrule
300 & 100 & E-CGL & 16.22 & 16.00 & 0.21 & 0.993 & 0.856 \\
300 & 100 & E-ACGL & 16.14 & 16.00 & 0.13 & 0.996 & 0.856 \\
300 & 200 & E-CGL & 16.22 & 16.00 & 0.22 & 0.993 & 0.859 \\
300 & 200 & E-ACGL & 16.09 & 16.00 & 0.08 & 0.997 & 0.859 \\
300 & 500 & E-CGL & 17.96 & 15.99 & 1.95 & 0.972 & 0.853 \\
300 & 500 & E-ACGL & 16.77 & 15.88 & 0.89 & 0.972 & 0.854 \\
1000 & 200 & E-CGL & 16.01 & 16.00 & 0.01 & 1.000 & 0.865 \\
1000 & 200 & E-ACGL & 16.01 & 16.00 & 0.01 & 1.000 & 0.865 \\
1000 & 500 & E-CGL & 16.05 & 16.00 & 0.05 & 0.998 & 0.867 \\
1000 & 500 & E-ACGL & 16.05 & 16.00 & 0.05 & 0.998 & 0.867 \\
\bottomrule
\end{tabular}
\end{table}

The effect of increasing $d$ was concentrated at $n=300$; at $n=1000$, both
centered group estimators remained close to the true support at $d=500$.

\section{Component-number diagnostics}
\label{sec:s-component-number}

The dense-model $K$ study uses $K\in\{2,\ldots,8\}$, $K^\ast=4$, $n=1000$,
$d=200$, and target $e_B=5\%$. Over 50 repetitions, the Akaike information
criterion (AIC) \citep{akaike1974aic} and independent test NLL selected $K=4$
in every repetition. The Bayesian information criterion (BIC)
\citep{schwarz1978bic} selected $K=4$ in 39/50 equal-concentration repetitions
and 0/50 heterogeneous repetitions. An EBIC-type criterion with index 0.5,
the risk inflation criterion (RIC), and its corrected form (RICc) selected
$K=2$ in all repetitions. A 10-repetition five-split
holdout diagnostic selected $K=4$ by minimum validation NLL in 10/10
repetitions under both concentration patterns; the one-standard-error (1-SE)
rule selected $K=4$
in 10/10 equal and 9/10 heterogeneous repetitions.

For a dense $K$-component fit, the diagnostic used
$\operatorname{df}_K=Kd+K-1$ and the following operational definitions:

\[
\begin{aligned}
\operatorname{AIC}(K)&=-2\ell_K+2\operatorname{df}_K,\\
\operatorname{BIC}(K)&=-2\ell_K+\log(n)\operatorname{df}_K,\\
\operatorname{RIC}(K)&=-2\ell_K+2\log(d)\operatorname{df}_K,\\
\operatorname{RICc}(K)&=-2\ell_K+2\{\log(d)+\log\log(d)\}\operatorname{df}_K,\\
\operatorname{EBIC\text{-}type}_{0.5}(K)&=-2\ell_K+
\{\log(n)+\log(d)\}\operatorname{df}_K.
\end{aligned}
\]

The EBIC-type criterion is an operational linear penalty inspired by the
large-model-space construction of \citet{chen2008ebic}; it is not their general
combinatorial EBIC expression.
The RIC and RICc labels follow \citet{foster1994ric} and
\citet{zhang2010ricc}, respectively. Here they refer to the recorded
high-dimensional penalty forms above and are used only as component-number
sensitivity diagnostics.

These findings motivate separation of component-number selection from support
selection. Predictive $K$ diagnostics are not presented as a theorem or as a
universally valid rule for high-dimensional mixtures.

\section{Rossi-style bridge diagnostic}
\label{sec:s-rossi-bridge}

A five-repetition bridge experiment follows the sparse-prototype setting of
\citet{rossi2022sparsevmf}: $n=200$, $d=100$, $K=4$, a target overlap of 5\%, and
10\% zero entries in the directional means. The mean achieved overlap was
4.70\%. Dense vMF, spherical $k$-means, and the Rossi-style M-L selector had mean
ARI 0.824, 0.750, and 0.814, respectively. E-CGL had mean ARI 0.591. Because
component-specific zero entries make the coordinate-union posterior-score contrast support
dense, this diagnostic targets sparse prototypes rather than the estimand of
E-CGL. It is retained as a literature bridge and not as primary evidence for
posterior-score contrast support recovery.

\section{Directional companion implementation}
\label{sec:s-target-diagnostic}

The centered mean-direction group lasso (M-CGL) is the directional companion.
Let $M\in\mathbb R^{K\times d}$ have row $k$ equal to $\mu_k^\top$ and let
$H_K=I_K-K^{-1}\mathbf 1_K\mathbf 1_K^\top$. The directional-contrast
companion solves
\[
\max_{\substack{
\pi\in[0,1]^K,\ \mathbf 1_K^\top\pi=1;\\
0\leq\kappa_k\leq\kappa_{\max},\
\ \lVert\mu_k\rVert_2=1,\ k=1,\ldots,K
}}
\left\{
\ell(\pi,\kappa,M)
-
\lambda_\mu\sum_{j=1}^{d}
\lVert H_KM_{\cdot j}\rVert_2
\right\}.
\]
The variable-splitting implementation uses the alternating direction method
of multipliers (ADMM) \citep{boyd2011admm}. It introduces an auxiliary matrix
$\Xi\in\mathbb R^{K\times d}$ subject to $\Xi=H_KM$,
alternates product-of-spheres gradient/retraction updates with centered group
thresholding and scaled-dual updates, and updates each $\kappa_k$ by numerical
inversion of the vMF mean-resultant relation. The Banerjee approximation
initializes and brackets the numerical inversion; it is not used as the final
M-CGL concentration estimate.

With fixed $M$, write

\[
\rho_k=\frac{r_k^\top\mu_k}{N_k},
\qquad
A_d(\kappa_k)=\rho_k.
\]

Values $\rho_k\leq0$ give $\widehat\kappa_k=0$, and positive values are clipped
below one. Here $A_d(0)=0$ by continuous extension. When the clipped value is
smaller than $A_d(\kappa_{\max})$, the root is found on
$[0,\kappa_{\max}]$; otherwise the cap is returned. The Banerjee initializer
\citep{banerjee2005vmf} is

\[
\kappa_{k,0}
=\frac{\rho_k(d-\rho_k^2)}{1-\rho_k^2},
\]

after which the numerical root, rather than this approximation, is retained.
If $\widehat\kappa_k=0$, the associated $\mu_k$ is not identified; interpretation
of directional-contrast support therefore presumes positive-concentration
components. All reported selected M-CGL fits satisfied this condition.

The same approximation has different roles elsewhere in the implementation.
The dense initializer and M-L use it as a capped concentration update, whereas
the penalized E-CGL block optimizes $\eta_k$ directly and does not invert
$A_d$; only its pooled path-scale calculation uses the approximation. For that
calculation, let

\[
\bar\rho_\bullet
=\frac{\lVert r_\bullet\rVert_2}{N_\bullet},
\qquad
\widetilde\rho_\bullet
=\min\{\max(\bar\rho_\bullet,10^{-10}),1-10^{-8}\},
\]

and define the final capped approximation

\[
\widehat\kappa_\bullet^B
=\min\left\{
\kappa_{\max},
\max\left\{
10^{-10},
\frac{\widetilde\rho_\bullet
(d-\widetilde\rho_\bullet^2)}
{1-\widetilde\rho_\bullet^2}
\right\}
\right\}.
\]

The implementation sets

\[
\eta^{\mathrm{common}}
=
\begin{cases}
\widehat\kappa_\bullet^B
r_\bullet/\lVert r_\bullet\rVert_2,
& \lVert r_\bullet\rVert_2>10^{-12},\\[3pt]
0, & \lVert r_\bullet\rVert_2\leq10^{-12}.
\end{cases}
\]

Thus, in the E-series path-scale calculation, the Banerjee approximation is
retained as the final pooled concentration rather than numerically inverted.
Consequently, the E-series endpoint is a KKT-motivated scale evaluated at the
approximate pooled fit, not the exact KKT threshold of the fully optimized
all-common fixed-responsibility subproblem.
Let $\varrho>0$ denote the augmented-Lagrangian penalty parameter. For ADMM
iteration $q$, the recorded residuals are

\[
R_{\mathrm{pri}}^{(q)}=H_KM^{(q)}-\Xi^{(q)},
\qquad
R_{\mathrm{dual}}^{(q)}
=\varrho H_K\{\Xi^{(q)}-\Xi^{(q-1)}\}
=\varrho\{\Xi^{(q)}-\Xi^{(q-1)}\}.
\]

The second equality holds because the $\Xi$ iterates remain column-centered.

For path construction, let $\widehat\Theta^{\mathrm D}_\mu$ be the dense
free-concentration pilot, let $\mathcal R^{\mathrm D}$ have row $k$ equal to
$(r_k^{\mathrm D})^\top$, define
$D_{\widehat\kappa^{\mathrm D}}
=\operatorname{diag}(\widehat\kappa_1^{\mathrm D},\ldots,
\widehat\kappa_K^{\mathrm D})$, and set
\[
B^{\mathrm D}
:=D_{\widehat\kappa^{\mathrm D}}\mathcal R^{\mathrm D},
\qquad
\lambda_{\mathrm{proxy},\mu}
:=\max_{1\le j\le d}
\lVert(H_KB^{\mathrm D})_{\cdot j}\rVert_2.
\]
The recorded M-CGL path is
\[
\Lambda_\mu
:=\{0\}\cup
\operatorname{Geom}
\left(
10^{-4}\lambda_{\mathrm{proxy},\mu},
1.25\lambda_{\mathrm{proxy},\mu};239
\right).
\]
This quantity is a dense-score scale proxy, not a global KKT threshold for the
nonconvex observed mixture objective.

The path implementation allows at most 500 ADMM iterations and uses scaled
primal and dual tolerances of $10^{-7}$; row-norm error on the product of
spheres is recorded separately. These are diagnostic controls for the
nonconvex companion, not a claim of general convex-ADMM convergence.

For a selected directional-contrast candidate $S$, the refit removes the
sparsity penalty and enforces
\[
j\notin S
\quad\Longrightarrow\quad
H_KM_{\cdot j}=0,
\]
while retaining component-specific concentrations. If $m=|S|$, the
generic nominal dimension is
\[
\operatorname{df}_{\mu,\mathrm{nom}}(S)
=
\begin{cases}
d+2K-2, & m=0,\\
d+(K-1)m+K-1, & m>0.
\end{cases}
\]
This is a practical BIC dimension under generic local-rank conditions, not an
effective degrees of freedom result. The M-CGL null model has a common
direction and free component-specific concentrations; it is not the pooled
single-vMF model.

The centered-direction penalty is bounded on the product of unit spheres and
therefore cannot prevent concentration divergence by itself. The directional
companion is consequently fitted with the same finite cap
$\kappa_{\max}=10^6$ as the E-series procedures.
This makes the computational objective well posed; it is not a sparsity
penalty or a claim that the unconstrained mixture likelihood is bounded.

Because the sphere constraint makes the split problem nonconvex, accepted
objective values, primal and dual residuals, and sphere feasibility are
reported as numerical checks. Global convergence is not claimed. The final
100-repetition estimand-separation results are reported once in the main
article; earlier path-development pilots are not repeated in this Supplementary Material.
Table~\ref{tab:s-mcgl-own-target} summarizes the four primary companion cells
using the directional target $S_\mu$ rather than the E-series target $S_\eta$.

\input{../simulations/tables/csda_supp_mcgl_own_target_260813.tex}

Table~\ref{tab:s-runtime-matched} reports the matched end-to-end runtime
diagnostic used to describe the relative computational burden.

The final timing panel was run sequentially in one R process on a Windows 10
Pro workstation with an Intel Core i7-11700K processor (8 physical and 16
logical cores) and 31.9 GiB of memory, using R 4.2.1 and Rcpp 1.0.12. The
\texttt{OMP\_NUM\_THREADS} and \texttt{MKL\_NUM\_THREADS} variables were unset,
so the underlying BLAS/OpenMP thread count was not explicitly pinned. One-time
Rcpp compilation was excluded. The table therefore documents implementation
burden on the recorded system, not hardware-independent algorithmic speed.

\input{../simulations/tables/csda_supp_runtime_diagnostics_260803.tex}

\clearpage
\section{Real-data stability and held-out diagnostics}
\label{sec:s-realdata-diagnostics}

The main article reports one fit to the full cleaned Classic3 payload.
This section retains the earlier five prespecified overlapping splits as
descriptive stability and held-out checks. The splitwise panel predates the
full-data sparse $k$-means addition and is not pooled with the full-data
estimates.

The Classic3 source counts were reconstructed from \texttt{Rcoclust} version
0.1.1. The source was
\href{https://cran.r-project.org/src/contrib/Archive/Rcoclust/Rcoclust_0.1.1.tar.gz}{the archived CRAN package},
and the data object was \texttt{data\_classic3.rda}. The archive version and
object path are fixed by the preparation script; labels from that object were
used only to set the benchmark resolution and evaluate fitted clusters.

The full-data payload was encoded with a maximum sequence length of 512 and
truncation enabled. The exact encoding identifiers were

\begingroup\footnotesize
\noindent Model:
\texttt{naver/splade-cocondenser-}\allowbreak\texttt{ensembledistil}\\
Revision:
\texttt{49cf4c7b0db5b870a401}\allowbreak\texttt{ddf5e2669993ef3699c7}.
\endgroup

Alphabetic tokens of length at least three with
document frequency at least two were eligible; the 2,000 largest unlabeled
full-payload variances were retained and rows were normalized to unit
Euclidean norm. Near-duplicate candidates required cosine similarity at least
0.90 and shingle Jaccard similarity at least 0.75. The final SHA-256 values
were

\begingroup\footnotesize
\noindent Vocabulary:\\
\texttt{9d4cb302c002212de31a6d5d8e48ad74}\allowbreak
\texttt{5b6390dcbb83bfe4260a963e3521c53c}\\
Payload:\\
\texttt{6f4e4842b6bc9e8f274714d4cbd9704}\allowbreak
\texttt{f1daacfbabf339a179037c09be7c096e3}.
\endgroup

The full-data ranking above applies only to the full-payload analysis in
the main article. For each repeated-holdout split, variance ranking was
recomputed from the training documents only; 2,000 Classic3 coordinates or
500 BBCSport coordinates were retained. The unchanged training-selected mask
was applied to the corresponding test documents, and the resulting training
and test rows were normalized separately. Thus the splitwise held-out results
are inductive with respect to feature ranking and do not use test labels or
test-document variances.

All reported NMI values were recomputed from the stored assignments using the
geometric normalization in Supplementary Section~\ref{sec:support-metrics}.
This harmonization did
not require model refitting or alter any displayed three-decimal full-data NMI.

The historical payload manifest did not retain exact Python, PyTorch, and
Transformers versions; none is inferred here.

\subsection{Sparse k-means restart check}

The full-data sparse $k$-means result was checked without changing its
label-blind permutation-gap rule. At the selected interior weight bound
3.869962, the recorded $n_{\mathrm{start}}=30$ fit had selected $q=985$,
ARI 0.137445, objective 35.63645, and cluster sizes 553, 2{,}546, and 784. Two
independent $n_{\mathrm{start}}=30$ batches at the same bound both selected
$q=1{,}024$, with ARI 0.137273, objective 35.63655, and cluster sizes 553, 2{,}545,
and 785. The adjacent lower and upper bounds gave ARI 0.136553 and 0.137445,
respectively. The low ARI was therefore reproducible across these independent
starts and adjacent tuning values under the SPLADE representation. This is a
local-optimum and implementation sanity check, not a claim about sparse
$k$-means on other representations.

\subsection{BBCSport data and analysis protocol}

BBCSport is the five-class document-clustering benchmark distributed by
\citet{greene2006diagonal}. The source archive contains 737 sports articles in
athletics, cricket, football, rugby, and tennis. Normalized-text screening
removed 13 exact duplicates; a subsequent audit removed 20 near-duplicates
using cosine similarity at least 0.90 followed by token-shingle Jaccard
similarity at least 0.75. The final analysis therefore used 704 documents,
with class counts 98, 116, 259, 139, and 92, respectively.

Five prespecified stratified 80/20 splits used seeds 20260730--20260734. Each
split contained 561 training and 143 test documents. SPLADE used the same
model and revision reported above. Within each split, coordinates were ranked
by their variance in the training documents only, the leading 500 were fixed,
and the same mask was applied to the test documents before separate rowwise
Euclidean normalization. Labels were used to form the stratified splits and
to evaluate fitted clusters, but not for coordinate ranking, initialization,
tuning, support selection, or parameter estimation. The class resolution was
fixed at $K=5$ to match the five supplied benchmark topics. Raw BBC articles
are not redistributed because the source content remains under BBC copyright.
Table~\ref{tab:s-realdata-full} summarizes clustering, support size, and
held-out fit across both datasets, and
Table~\ref{tab:s-realdata-e-series-path-audit} reports the Classic3 E-series
historical path sensitivity on the same splits.

\input{../simulations/tables/csda_supp_realdata_five_split_results_floor1e4_260814.tex}

\input{../simulations/tables/csda_supp_realdata_e_series_split_path_audit_260808.tex}

\subsection{Paired sparse--dense comparisons}

Table~\ref{tab:s-realdata-paired} reports paired differences between each
sparse method and the dense model having the same concentration structure.
M-L is paired with dense shared-$\kappa$ vMF; M-CGL and the E-series methods are paired
with dense free-$\kappa_k$ vMF. Positive NLL differences indicate worse
held-out density.

\input{../simulations/tables/csda_supp_realdata_paired_floor1e4_260814.tex}

\FloatBarrier

Because the holdouts overlap, these counts are descriptive. In particular,
the table is not a noninferiority analysis. Selected $q$ is also not a common
estimand across prototype, directional-contrast, and posterior-score
contrast methods.

The all-common candidate was also evaluated directly for the full-data
Classic3 E-series fits and for each BBCSport training split. It is a
single-vMF candidate with nominal dimension $d$; no penalized path was rerun
for this audit. Table~\ref{tab:s-realdata-empty} shows that it did not minimize BIC
in any of the 12 comparisons.

\input{../simulations/tables/csda_supp_realdata_empty_support_floor1e4_260812.tex}

Fig.~\ref{fig:s-realdata-splits} compares held-out clustering accuracy and
method-specific coordinate retention across the five splits.

\begin{figure}[H]
\centering
\includegraphics[width=\textwidth]{figure_06.png}
\caption{Clustering accuracy and coordinate retention across the five
prespecified real-data splits. Coordinate retention is defined relative to
each method-specific support target. E-CGL and E-ACGL use the final single
240-point KKT-motivated path in both datasets; M-CGL uses its 240-point
dense-score-proxy path, and M-L retains its event-driven path.}
\label{fig:s-realdata-splits}
\end{figure}

\clearpage
\subsection{Classic3 support and token stability}

For E-CGL and E-ACGL, the conditional support Jaccard index across the prespecified
training vocabularies was 0.933 under the final method-specific single-path
rule. Coordinates absent
from one of a pair of training vocabularies were excluded from that pair's
denominator. The following positive-contrast tokens were available and
selected by E-CGL in all five splits. Tables~\ref{tab:s-realdata-support-stability}
and \ref{tab:s-realdata-figure4-token-frequency} summarize support overlap and
the splitwise selection frequencies of the displayed contrast tokens.

Grouped by the post-fit component labels, the repeated leading positive tokens
were \emph{library}, \emph{information}, \emph{librarian}, \emph{libraries},
and \emph{retrieval} for CISI; \emph{flow}, \emph{mach}, \emph{pressure},
\emph{heat}, and \emph{theory} for CRAN; and \emph{tumor}, \emph{inhibitor},
\emph{rat}, \emph{dose}, and \emph{cancer} for MED.

Labels were used only after fitting to name components and display the
contrasts. An excluded E-CGL coordinate was constrained to have no centered
natural-parameter contrast; its common baseline was retained.

\FloatBarrier

\input{../simulations/tables/csda_supp_realdata_support_stability_floor1e4_260812.tex}

\input{../simulations/tables/csda_supp_figure4_panelA_token_stability_floor1e4_260812.tex}

\FloatBarrier

\clearpage
\subsection{BBCSport M-L path sensitivity}

The primary M-L path allowed at most 600 path points, including the
zero-penalty fit, and imposed a minimum relative penalty increment of 0.02. A
sensitivity run allowed at most 1,200 path points, including the zero-penalty
fit, and reduced the minimum increment to 0.005. Results are summarized in
Table~\ref{tab:s-bbcsport-ml-path}.

\begin{table}[H]
\centering
\small
\caption{Path-resolution sensitivity of M-L in BBCSport. Path settings and
mean diagnostics are reported over the five prespecified splits.}
\label{tab:s-bbcsport-ml-path}
\begin{tabular}{lrr}
\toprule
Quantity & Primary path & Denser path \\
\midrule
Maximum path points & 600 & 1,200 \\
Minimum relative increment & 0.020 & 0.005 \\
Mean evaluated path rows & 331.8 & 771.2 \\
Mean selected $q$ & 498.2 & 498.4 \\
Mean M-L runtime ratio & 1.00 & 1.76 \\
\bottomrule
\end{tabular}
\end{table}

Selected $q$ differed by at most one coordinate at the split level, and
held-out ARI/NMI were unchanged in all five splits. The final evaluated
support had BIC at least 19,497.94 larger than the selected support. Both
paths nevertheless stopped after a later stronger-penalty fit failed. The
near-dense M-L selection is therefore stable to the examined path resolution,
but no claim is made about unevaluated candidates beyond the numerical
failure.

\subsection{Numerical audit}

Table~\ref{tab:s-realdata-numerics} summarizes completion, convergence,
boundary, and path-stop diagnostics for the prespecified split analyses.

\begin{table}[H]
\centering
\small
\caption{Convergence and boundary diagnostics for the methods included in the
final prespecified split-performance tables.}
\label{tab:s-realdata-numerics}
\resizebox{\textwidth}{!}{%
\begin{tabular}{lrrrrr}
\toprule
Dataset & Completed & Errors & Nonconverged & Boundary selected
  & Path fit stops \\
\midrule
Classic3 & 35/35 & 0 & 0 & 0 & 0 \\
BBCSport & 35/35 & 0 & 0 & 0 & 5 \\
\bottomrule
\end{tabular}
}
\end{table}

Classic3 M-CGL emitted high-order Bessel precision warnings. Over the
selected-fit range, $d=2{,}000$ and $\kappa\in[607,818]$, the production
Bessel-ratio calculation was compared with a reference based on the
second-order uniform large-order expansion of the modified Bessel function
\citep[Sec.~10.41]{olver2010nist}. With $\nu=d/2-1$, $z=\kappa/\nu$, and
$p={(1+z^2)}^{-1/2}$, the reference retained

\[
I_\nu(\nu z)
\simeq
\frac{\exp(\nu\zeta)}{\sqrt{2\pi\nu}{(1+z^2)}^{1/4}}
\left\{1+\frac{U_1(p)}{\nu}+\frac{U_2(p)}{\nu^2}\right\},
\]

where

\[
\begin{aligned}
\zeta
&=\sqrt{1+z^2}+\log\frac{z}{1+\sqrt{1+z^2}},\\
U_1(p)
&=\frac{3p-5p^3}{24},
&
U_2(p)
&=\frac{81p^2-462p^4+385p^6}{1152}.
\end{aligned}
\]

Writing $C_d^{(2)}$ for the vMF normalizer induced by this second-order
approximation, the reference ratio was evaluated by the centered finite
difference

\[
A_d^{\mathrm{ref}}(\kappa)
=-
\frac{\log C_d^{(2)}(\kappa+h)-\log C_d^{(2)}(\kappa-h)}{2h},
\qquad
h=10^{-5}\max(1,\kappa).
\]

The maximum absolute and relative differences were
$8.983\times10^{-11}$ and $2.850\times10^{-10}$, respectively. The observed
range passed the prespecified relative tolerance of $10^{-6}$; this is not a
global approximation-error bound.

A prespecified concentration-guard audit then refit E-CGL and E-ACGL on all five
Classic3 and five BBCSport splits under the common
$\texttt{kappa\_cap}=10^6$ rule. Across the 20 dataset--split--method
comparisons, selected supports and held-out assignments agreed in 20/20
cases. The maximum absolute changes in train log-likelihood, test
log-likelihood, and the fitted natural parameters were all zero at the stored
precision. The largest path and refit ratios
$\max_k\widehat\kappa_k/\texttt{kappa\_cap}$ were 0.002525 and 0.002537,
respectively, and all refits passed the 0.99 eligibility guard. This is a
splitwise equivalence audit, not a bound over every possible degeneracy
configuration.

\section{CSTR Rossi-style reproduction bridge}
\label{sec:s-cstr-bridge}

The Computer Science Technical Reports (CSTR) collection contains 475
computer-science technical-report abstracts represented by
1,000 binary word coordinates and four reference subject areas, following the
benchmark used by \citet{rossi2022sparsevmf}. The implementation reproduced
the published clustering levels for the dense shared-$\kappa$ mixture and the
Rossi M-L selector, as summarized in Table~\ref{tab:s-cstr-bridge}.

\begin{table}[H]
\centering
\small
\caption{CSTR comparison with the published Rossi-style analysis.
Implementation results are mean (SD) over 50 runs.}
\label{tab:s-cstr-bridge}
\begin{tabular}{lrr}
\toprule
Method & Published ARI & Implementation ARI \\
\midrule
Dense vMF, shared $\kappa$ & 0.804 & 0.8023 (0.0087) \\
Rossi M-L, BIC & 0.808 & 0.8083 (0.0079) \\
\bottomrule
\end{tabular}
\end{table}

The earlier centered-$\eta$ CSTR diagnostic used BIC before refit and is not
combined numerically with the prespecified BIC-after-refit analysis. CSTR is used only to
check comparability of the Rossi-style implementation.

\bibliographystyle{elsarticle-harv}
\bibliography{references}

\end{document}
